\section{Arbitrary coordinate domains and local graph lifts}
\label{sec:coordinate-domains}
The endpoint argument uses membership, not interval geometry. We first
strengthen the information available about each original coordinate, then
show that adding arbitrarily many auxiliaries depending on that coordinate
does not cost degree in this particular moment construction.

\subsection{All valid local polynomial information}
Let $S=\prod_i S_i$, where each $S_i$ is an arbitrary nonempty subset of
$[0,1]$. Closedness, connectedness, semialgebraicity, and a finite description
are not required. Define the order-$r$ product-domain oracle by normalized
functionals on $\R[x]_{\le2r}$, all balance products
\eqref{eq:preorder-equality}, and all inequalities
\begin{equation}\label{eq:product-domain-oracle}
 L\left[\left(\prod_{j=1}^a g_j\right)P^2\right]\ge0,
 \qquad\sum_j\deg g_j+2\deg P\le2r,
\end{equation}
where each $g_j$ is any univariate polynomial nonnegative on its
coordinate set $S_i$. The empty product is allowed, factors may repeat,
and $P$ may couple every coordinate. We also grant
\begin{equation}\label{eq:product-domain-equalities}
 L[h v]=0\qquad(\deg h+\deg v\le2r)
\end{equation}
for every univariate polynomial $h$ vanishing on its $S_i$ and every
polynomial multiplier $v$ in the original coordinates. In particular this
model grants all local information through degree, not merely a chosen
list of input generators. Its representation need not be finite or
computationally efficient. The node bound is the infimum of $L[F]$ or
$L[F_d]$, as appropriate. A certified cover has original feasible set
$\mathcal F_{n,k}$, original dimension $n$, and target
$1/4-\epsilon$ or $F_d^*-\epsilon$ respectively.

\begin{theorem}[Product-domain certificates]\label{thm:product-domain-cover}
Under the respective parameter assumptions of
Theorem~\ref{thm:preordering-cover} and
Corollary~\ref{cor:unique-cardinality}, the cover bounds
$\mathcal N(q_r)$ and $\mathcal N(q_{r,\eta})$ hold for the
product-domain oracle~\eqref{eq:product-domain-oracle}--\eqref{eq:product-domain-equalities}.
\end{theorem}
\begin{proof}
For a domain containing a witness, use the exclusions
\eqref{eq:endpoint-exclusions}. If $|R|<q_r$, or $|R|<q_{r,\eta}$ in
the perturbed case, fix $x_R=w_R$ and use $E_{s,t}$ on $U$ exactly as
in Theorem~\ref{thm:preordering-cover}. Its parameter and balance
calculations are unchanged.

A local inequality on a restricted coordinate becomes a nonnegative
constant because $w_i\in S_i$. For $i\in U$, both endpoints belong to
$S_i$, and its Boolean reduction is
\begin{equation}\label{eq:local-endpoint-interpolation}
 \operatorname{red}g(u_i)=g(0)(1-u_i)+g(1)u_i,
 \qquad g(0),g(1)\ge0.
\end{equation}
Group factors belonging to the same coordinate and expand their product
in endpoint indicators. After removing constant factors, let $J$ be the
coordinates with nonconstant reduced factors. The expansion has the form
\[
 \sum_{\sigma\in\{0,1\}^J} c_\sigma
       \prod_{i\in J:\sigma_i=1}u_i
       \prod_{i\in J:\sigma_i=0}(1-u_i),\qquad c_\sigma\ge0.
\]
Every such coordinate requires at least one original degree, so
$|J|\le\sum_j\deg g_j$. If $\widehat P$ is the substituted square
polynomial, $\deg\widehat P\le\deg P$ and
$|J|+2\deg\widehat P\le2r$.
Lemma~\ref{lem:assignment-localizer} gives nonnegativity of each
indicator times $\widehat P^2$, hence every inequality in
\eqref{eq:product-domain-oracle}. This includes repeated factors and
squares coupling restricted and unrestricted coordinates.

A locally vanishing $h$ becomes zero at $w_i$ on $R$, or has zero
values at both endpoints on $U$. Its Boolean reduction is zero. This
remains true after multiplication by every allowed, possibly globally
coupled, $v$, proving \eqref{eq:product-domain-equalities}.
The objective value and perturbation estimate are exactly
\eqref{eq:restricted-penalty} and Corollary~\ref{cor:unique-cardinality}.
Thus certification forces the same restriction count, and
Lemma~\ref{lem:endpoint-count} finishes the proof.
\end{proof}
The constructed first moment need not itself belong to a disconnected
$S_i$; it is a moment satisfying the stipulated oracle. The proof uses only
endpoint and witness membership. It invokes no compact-domain hull theorem
and takes no closure of $S_i$.
A branch on $\psi_i(x_i)\le\beta$ versus $\psi_i(x_i)\ge\beta$ changes
only $S_i$, for any real-valued univariate function $\psi_i$. In particular,
branching on $x_i(1-x_i)$, with its disconnected endpoint-side region,
is covered. For an objective-based reduction $S_i\to T_i\subset S_i$,
the removed product region has coordinate set $S_i\setminus T_i$.
If nonempty and certified at the target, it can be charged directly under
Proposition~\ref{prop:charged-tolerance}, without a closure convention.

\subsection{Endpoint interpolation in lifted coordinates}
For each original coordinate $i$, introduce finitely many auxiliaries
\begin{equation}\label{eq:local-graph-map}
 y_{i,j}=\psi_{i,j}(x_i),\qquad j=1,\ldots,a_i,
 \qquad \psi_{i,j}:[0,1]\to\R.
\end{equation}
Every function is defined and finite on the entire interval. Neither
continuity nor polynomiality is required. Write
$\Gamma_i(v)=(v,\psi_{i,1}(v),\ldots,\psi_{i,a_i}(v))$.
A node restricts each local graph independently to
$\Gamma_i(S_i)$ for arbitrary nonempty $S_i\subseteq[0,1]$.
Thus restrictions in the lifted variables are measured by their exact
preimages. Original coordinates $x_i$ are retained and the original
balances remain affine. The number of auxiliaries is finite for each
instance but otherwise unrestricted.

The lifted order-$r$ oracle acts on polynomials in all retained and
auxiliary coordinates of total \emph{lifted} degree at most $2r$. It
imposes normalization; all multiples of the original affine balance
through that degree; every local polynomial equality vanishing on
$\Gamma_i(S_i)$ times every allowed global polynomial multiplier; and
\begin{equation}\label{eq:lifted-local-oracle}
 L\left[\left(\prod_j g_j\right)P^2\right]\ge0,
 \qquad\sum_j\deg g_j+2\deg P\le2r,
\end{equation}
where each $g_j$ is a polynomial in one local graph block, nonnegative
on that block's restricted graph. Degrees in this display are measured
before substitution. Squares can couple every lifted block, and factors
may repeat. Graph equations are included whenever their lifted degrees
allow them. Arbitrary coupled valid inequalities are not supplied.

Let $\widetilde F(x,y)$ be a polynomial of lifted degree at most $2r$
that agrees exactly with $F(x)$, or with $F_d(x)$, on the \emph{full}
product graph $\prod_i\Gamma_i([0,1])$. Agreement merely on the
balance-feasible graph is not the stated assumption. The node bound is
$\inf L[\widetilde F]$. A graph-region cover means that its preimages
cover the original compact set $\mathcal F_{n,k}$. Full-graph objective
agreement preserves its optimum and an attaining optimizer, even if the
lifted graph itself is noncompact because some $\psi_{i,j}$ is
noncontinuous or unbounded near a point. Targets are still
$T_*=1/4-\epsilon$ or $T_*=F_d^*-\epsilon$.

\begin{theorem}[No order loss for coordinatewise graph lifts]
\label{thm:local-graph-lift}
For the lifted model just defined, under the respective assumptions of
Theorem~\ref{thm:preordering-cover} and
Corollary~\ref{cor:unique-cardinality}, every certified cover has at
least $\mathcal N(q_r)$ or $\mathcal N(q_{r,\eta})$ regions.
The parameter is the lifted order $r$ itself, independent of the number
or polynomial degrees of the univariate auxiliaries. The dimension in
the bound is the number $n$ of original coordinates.
\end{theorem}
\begin{proof}
Fix a node containing a graph witness $\Gamma(w)$ and form $A,D,R,U$
from its preimage sets. Suppose $|R|$ is below the appropriate threshold.
The cardinality parameters $s,t$ satisfy the same hypotheses as before.
Define an algebraic substitution $\Phi$ by
\begin{align}
 \Phi(x_i)&=w_i,&
 \Phi(y_{i,j})&=\psi_{i,j}(w_i) &&(i\in R),\notag\\
 \Phi(x_i)&=u_i,&
 \Phi(y_{i,j})&=\psi_{i,j}(0)
       +(\psi_{i,j}(1)-\psi_{i,j}(0))u_i &&(i\in U).
 \label{eq:graph-interpolation-map}
\end{align}
Every image is affine or constant, so
$\deg\Phi(P)\le\deg P$ for every lifted polynomial. Define
\[
 L[P]=E_{s,t}[\Phi(P)]\qquad(\deg P\le2r).
\]
This is well-defined and normalized. The original affine balance maps
to $\sum_{i\in U}u_i-t$, and any allowed multiplier maps to degree at
most $2r-1$. Lemma~\ref{lem:fractional-positive} proves all balance
products.

A local valid inequality $g$ maps to a nonnegative constant on $R$.
For $i\in U$, the two values of $\Phi(g)$ are
$g(\Gamma_i(0))$ and $g(\Gamma_i(1))$. Both are nonnegative because
both endpoint graph tuples belong to the restricted local graph.
Thus its Boolean reduction has the nonnegative endpoint form
\eqref{eq:local-endpoint-interpolation}. Grouping and expanding all
local factors gives nonnegative combinations of assignment indicators
on at most $\sum_j\deg g_j$ distinct Boolean coordinates. Moreover
$\deg\Phi(P)\le\deg P$, so each indicator times $\Phi(P)^2$ lies
in the degree range of Lemma~\ref{lem:assignment-localizer}.
This proves every inequality~\eqref{eq:lifted-local-oracle}, including
squares coupling all original and auxiliary variables.

For a local equality $h$, $\Phi(h)$ vanishes at the deterministic
witness tuple on $R$ or at both endpoint graph tuples on $U$. Its
Boolean reduction is zero. For any allowed multiplier $v$,
$\operatorname{red}(\Phi(h)\Phi(v))=0$, and the degree bound ensures
$L[hv]$ is defined. Thus every granted local equality product holds.
The same reasoning would also respect any polynomial identity valid on
the full product graph, even a coupled identity, and its allowed
multiples: at every Boolean assignment the image is an actual graph
point. This observation does not grant coupled inequalities that use
the balances or objective cutoff.

At each Boolean assignment of $u_U$, \eqref{eq:graph-interpolation-map}
is an actual full-graph tuple with $x_R=w_R$. Therefore
$\Phi(\widetilde F)$ agrees at every Boolean point with the polynomial
obtained by substituting $x_R=w_R,x_U=u$ in $F$ or $F_d$. Two
polynomials agreeing on the Boolean cube have the same multilinear
reduction: successively evaluating at zero and one in each coordinate
proves uniqueness of that representation. Both polynomials have degree
at most $2r$. Their functional values are consequently identical.
The penalty is $|M\cap R|p(1-p)$ and any perturbation contributes at
most $\eta$. The strict below-target estimate and the endpoint count
from the preceding proofs finish the theorem.
\end{proof}

The substitution $\Phi$ constructs a moment functional. It is not an
identity $\psi_i(x_i)=\psi_i(0)+(\psi_i(1)-\psi_i(0))x_i$ for continuous
$x_i$. The proof checks true graph identities only at deterministic or
Boolean endpoint tuples. For example, it covers a square-root auxiliary
$y_i=\sqrt{x_i}$ with its polynomial identity $y_i^2=x_i$, and a
penalty auxiliary $y_i=x_i(1-x_i)$ with linear objective $\sum_i y_i$.
Both maps are finite on all of $[0,1]$ and their objective agreement, when
used to encode the intended objective, is checked on the full graph.
It also covers finite-valued nonpolynomial maps that have few polynomial
identities; the oracle still grants every local polynomial inequality
valid on the restricted graph.

The theorem requires independent local restrictions and retained original
coordinates. Eliminating $x_i$ through a nonlinear inverse, branching on
functions of several original coordinates, granting arbitrary coupled
valid cuts, or replacing the oracle by the exact hull of the coupled
feasible graph changes the model. Arbitrarily many auxiliaries also mean
that an exponential bound in original dimension need not be exponential
in lifted dimension.

\begin{proposition}[Literal polynomial substitution as a coarser bound]
\label{prop:literal-lift}
If all $\psi_{i,j}$ are polynomials of degree at most an integer $D\ge1$,
with the originals counted as degree one, literal graph substitution
also proves the lifted cover bounds with effective order $rD$ in place
of $r$. This is a valid weaker conclusion of the same type.
\end{proposition}
\begin{proof}
Let $\Psi$ replace each auxiliary by its defining polynomial and fix
no coordinates yet. Then $\deg\Psi(P)\le D\deg P$.
Take the product-domain functional of order $rD$ from
Theorem~\ref{thm:product-domain-cover} and set $\widetilde L[P]=L[\Psi(P)]$.
A lifted local inequality pulls back to a univariate polynomial valid on
$S_i$, and
\[
 \sum_j\deg\Psi(g_j)+2\deg\Psi(P)
 \le D\left(\sum_j\deg g_j+2\deg P\right)\le2rD.
\]
Thus every localizer holds; local equalities pull back to equalities
vanishing on $S_i$, and graph identities vanish identically.
The original balance has degree one after substitution, so for a lifted
multiplier $v$ of degree at most $2r-1$,
\[
 \deg(e\Psi(v))\le1+D(2r-1)\le2rD.
\]
All equality products are therefore available. Objective agreement on
the full graph means the substituted objective equals $F$ or $F_d$
on $[0,1]^n$. Since both are polynomials and the cube has interior, this
is a polynomial identity. Hence its evaluated objective is the required
one. The product-domain count at order $rD$ applies.
\end{proof}
For quadratic penalty auxiliaries the literal argument uses
$\min\{k-4r+2,z-4r+2,m(1/2-2\epsilon)\}$.
Theorem~\ref{thm:local-graph-lift} improves this to $q_r$ by using
endpoint interpolation instead of literal substitution. It is this
specific construction, not invariance of polynomial hierarchies under
arbitrary nonlinear reformulations, that removes the degree loss.
